% Options for packages loaded elsewhere
\PassOptionsToPackage{unicode}{hyperref}
\PassOptionsToPackage{hyphens}{url}
\documentclass[
  12pt,
]{article}
\usepackage{xcolor}
\usepackage{amsmath,amssymb}
\setcounter{secnumdepth}{-\maxdimen} % remove section numbering
\usepackage{iftex}
\ifPDFTeX
  \usepackage[T1]{fontenc}
  \usepackage[utf8]{inputenc}
  \usepackage{textcomp} % provide euro and other symbols
\else % if luatex or xetex
  \usepackage{unicode-math} % this also loads fontspec
  \defaultfontfeatures{Scale=MatchLowercase}
  \defaultfontfeatures[\rmfamily]{Ligatures=TeX,Scale=1}
\fi
\usepackage{lmodern}
\ifPDFTeX\else
  % xetex/luatex font selection
\fi
% Use upquote if available, for straight quotes in verbatim environments
\IfFileExists{upquote.sty}{\usepackage{upquote}}{}
\IfFileExists{microtype.sty}{% use microtype if available
  \usepackage[]{microtype}
  \UseMicrotypeSet[protrusion]{basicmath} % disable protrusion for tt fonts
}{}
\makeatletter
\@ifundefined{KOMAClassName}{% if non-KOMA class
  \IfFileExists{parskip.sty}{%
    \usepackage{parskip}
  }{% else
    \setlength{\parindent}{0pt}
    \setlength{\parskip}{6pt plus 2pt minus 1pt}}
}{% if KOMA class
  \KOMAoptions{parskip=half}}
\makeatother
\usepackage{longtable,booktabs,array}
\usepackage{calc} % for calculating minipage widths
% Correct order of tables after \paragraph or \subparagraph
\usepackage{etoolbox}
\makeatletter
\patchcmd\longtable{\par}{\if@noskipsec\mbox{}\fi\par}{}{}
\makeatother
% Allow footnotes in longtable head/foot
\IfFileExists{footnotehyper.sty}{\usepackage{footnotehyper}}{\usepackage{footnote}}
\makesavenoteenv{longtable}
\setlength{\emergencystretch}{3em} % prevent overfull lines
\providecommand{\tightlist}{%
  \setlength{\itemsep}{0pt}\setlength{\parskip}{0pt}}
\usepackage{bookmark}
\IfFileExists{xurl.sty}{\usepackage{xurl}}{} % add URL line breaks if available
\urlstyle{same}
\hypersetup{
  hidelinks,
  pdfcreator={LaTeX via pandoc}}

\author{SCX}
\date{2026}

\begin{document}

\section{Situs (Physical Positional Encoding): Physical Reality Verification and Applicability Scenario Analysis}\label{situs-physical-reality-verification}

\begin{center}\rule{0.5\linewidth}{0.5pt}\end{center}

\textbf{Document Status}: Honest critique --- no euphemisms\\
\textbf{Prerequisites}: \texttt{ppe\_rigorous\_derivation.md} (rigorous mathematical derivation), \texttt{multi\_head\_spring\_and\_positional\_encoding\_analysis.md} (CC audit report)\\
\textbf{Methodology}: Real physical data + physical correspondence of mathematical theorems + scenario-by-scenario argumentation\\
\textbf{Core Criterion}: Theorem 2.2.1 --- $\delta_s^{\text{PE}} > 0 \iff I(Y; P \mid S) > 0$, i.e., physical position must carry additional information about the label given the state atom condition.

\begin{center}\rule{0.5\linewidth}{0.5pt}\end{center}

\section{Part 1: Physical Reality Checks}\label{part-1-physical-reality-checks}

\begin{center}\rule{0.5\linewidth}{0.5pt}\end{center}

\subsection{1.1 Protein Sequence Positional Encoding}\label{protein-sequence-positional-encoding}

\subsubsection{1.1.1 Same Amino Acid, Different Position: Lys-37 (Active Site) vs Lys-289 (Surface Loop)}\label{lys-37-vs-lys-289}

\textbf{Biological facts (indisputable)}:

The same lysine residue at different positions in a protein has completely different functional roles:
- \textbf{Active site Lys-37}: Typical catalytic residue, side-chain pKa can deviate 2--4 units from solution value (from $\sim$10.5 down to 6.0--8.0), participating in acid-base or covalent catalysis. Examples: Lys-115 in acetoacetate decarboxylase with pKa $\approx$ 6.0 [Bartlett et al., \emph{J. Mol. Biol.} 324, 105 (2002)]; Lys-41 in ribonuclease A with pKa $\approx$ 8.8 [Raines, \emph{Chem. Rev.} 98, 1045 (1998)]. Highly conserved, extremely sensitive to mutation ($\Delta\Delta$G $>$ 3 kcal/mol common).
- \textbf{Surface loop Lys-289}: Solvent-exposed, pKa close to solution value (10.4--10.6), functionally involved in non-specific electrostatic interactions or protein-protein recognition interfaces, low conservation, tolerant to mutation ($\Delta\Delta$G typically $<$ 1 kcal/mol).

\textbf{Physical judgment}:
\[I(Y; P \mid S) > 0 \quad \text{--- given ``amino acid type = Lys'', residue position carries extensive functional information}\]

\begin{longtable}[]{@{}
  >{\raggedright\arraybackslash}p{(\linewidth - 6\tabcolsep) * \real{0.1311}}
  >{\raggedright\arraybackslash}p{(\linewidth - 6\tabcolsep) * \real{0.3443}}
  >{\raggedright\arraybackslash}p{(\linewidth - 6\tabcolsep) * \real{0.3607}}
  >{\raggedright\arraybackslash}p{(\linewidth - 6\tabcolsep) * \real{0.1639}}@{}}
\toprule\noalign{}
\begin{minipage}[b]{\linewidth}\raggedright
Physical Quantity
\end{minipage} & \begin{minipage}[b]{\linewidth}\raggedright
Lys-37 (Active Site)
\end{minipage} & \begin{minipage}[b]{\linewidth}\raggedright
Lys-289 (Surface Loop)
\end{minipage} & \begin{minipage}[b]{\linewidth}\raggedright
Source of Difference
\end{minipage} \\
\midrule\noalign{}
\endhead
\bottomrule\noalign{}
\endlastfoot
Side-chain pKa & 6.0--8.0 & 10.4--10.6 & Local electrostatic environment (buried vs solvent-exposed) \\
SASA (Solvent Accessible Surface Area) & $<$ 10 \AA$^2$ & $>$ 100 \AA$^2$ & 3D fold position \\
Evolutionary conservation (BLOSUM62 position entropy) & $<$ 0.2 bits & $>$ 0.8 bits & Functional constraint strength \\
Mutation $\Delta\Delta$G & $>$ 3 kcal/mol & $<$ 1 kcal/mol & Folding stability contribution \\
\end{longtable}

\textbf{BUT --- Critical Limitation}:

Situs's sine encoding encodes \textbf{sequence position} (1D residue number 37 vs 289), not \textbf{3D spatial position} (active site pocket vs surface). This is a fundamental flaw of Situs in protein applications:

\begin{itemize}
\tightlist
\item
  Sequence position $\to$ 3D position mapping is \textbf{not a functional relationship}. In homologous proteins, sequence position 37 may be the active site in homolog A but a loop in homolog B (insertions/deletions cause sequence shifts).
\item
  For cases without 3D structure (most protein prediction tasks), Situs can only use sequence position. Sequence position is a \textbf{weak proxy} for 3D functional position.
\item
  Only when we have 3D structure and apply Situs to 3D coordinate encoding can the above functional differences be effectively captured.
\end{itemize}

\textbf{Honest conclusion}: 1D sequence sine encoding for distinguishing Lys-37 active site vs Lys-289 surface loop is \textbf{informative but the signal is weak}. The correlation between sequence position and 3D functional position is limited by fold conservation. $I(Y; P_{\text{1D}} \mid S) > 0$ but far smaller than $I(Y; P_{\text{3D}} \mid S)$.

\begin{center}\rule{0.5\linewidth}{0.5pt}\end{center}

\subsubsection{1.1.2 The 3.6 Residues/Turn Periodicity of $\alpha$-Helices}\label{alpha-helix-periodicity}

\textbf{Physical facts}:

Pauling's classic $\alpha$-helix [Pauling et al., \emph{PNAS} 37, 205 (1951)]:
- 3.6 residues per turn (exact value = 3.6, or 18 residues/5 turns).
- Pitch = 5.4 \AA\ (1.5 \AA\ rise per residue).
- Hydrogen bond pattern: residue $i$'s C=O $\cdots$ H--N of residue $i+4$.
- This means residues $i$ and $i+4$ are spatially adjacent ($\sim$3.1 \AA), while $i$ and $i+3$ (or $i+5$) are on opposite faces of the helix.

\textbf{Can sine encoding capture this?}

For sine encoding $\sin(2\pi p / \lambda_j)$:
- If wavelength $\lambda_j = 3.6$: $\sin(2\pi p/3.6)$ indeed has exact period 3.6 residues.
- Residues $i$ and $i+3.6$ get the same encoding value $\to$ residues $i$ and $i+4$ (distance 4, period difference 0.4/3.6 $\approx$ 11\%) have close but not identical encodings.

\textbf{However --- standard Transformer encoding cannot do this}:

The wavelength set of standard sine encoding [Vaswani et al., \emph{NeurIPS} 2017]:
\[\lambda_j = 2\pi \cdot 10000^{2j/d}\]

When $d = 512$, wavelengths range from approximately $2\pi$ to $2\pi \cdot 10000$ ($\approx$ 6.28 to $\approx$ 62832). For typical protein length $L \sim 300$:
- Shortest wavelength ($j = d/2-1 = 255$): $\lambda_{255} = 2\pi \cdot 10000^{510/512} \approx 2\pi \cdot 10000^{0.996} \approx 6.28 \times 9550 \approx 60,000$ --- far larger than 3.6!
- Not a single $\lambda_j$ is close to 3.6.

\textbf{Can the physically optimal spectrum (Theorem 1.2.1) solve this?}

According to Theorem 1.2.1, the optimal frequency spectrum requires knowing the physical correlation length $\xi$. For $\alpha$-helices, $\xi \approx 3.6$ (helix period). If we design the encoding targeting $\xi = 3.6$:
- $\lambda_{\min} = 2\pi\xi \cdot \cot(\pi(2(d/2-1)+1)/2d) \approx 2\pi \cdot 3.6 \cdot \cot(\pi/2 - \pi/2d) \approx 22.6 \cdot (2d/\pi) \approx 14.4 d$
- For $d = 64$, $\lambda_{\min} \approx 920$ --- still far larger than 3.6!

\textbf{Fundamental problem}: The ``resolution'' of sine encoding is limited by $d$. To resolve a 3.6-residue period on a sequence of length $L = 300$, we need $\lambda_{\min} \lesssim 3.6$, which requires $d \gtrsim 300/3.6 \times 2 = 167$ dimensions (at least 2 sample points per period, Nyquist). For typical $d = 64$ or $d = 128$, the 3.6-residue periodicity \textbf{cannot be captured in theory}.

\textbf{Honest conclusion}: For practically used encoding dimensions ($d \leq 128$), sine encoding \textbf{cannot} reliably capture the 3.6 residues/turn periodicity of $\alpha$-helices. This is not a defect of the encoding form but a fundamental dimension-resolution tradeoff (Corollary 1.2.1). If one needs to capture helical periodicity in sequence position encoding, one should explicitly add a frequency component with $\lambda \approx 3.6$, or use learnable periodic features.

\begin{center}\rule{0.5\linewidth}{0.5pt}\end{center}

\subsubsection{1.1.3 Degradation of Situs in Intrinsically Disordered Regions (IDRs)}\label{idr-degradation}

\textbf{Physical facts}:

Intrinsically disordered regions (IDRs) constitute 30--40\% of the eukaryotic proteome [Dunker et al., \emph{FEBS J.} 272, 5129 (2005)]. Key features:
- \textbf{Do not fold} into a single stable 3D structure under physiological conditions.
- Function arises from \textbf{amino acid composition} (e.g., charge patterns, proline content) rather than precise folding.
- Sequence position-function relationship is \textbf{extremely weak}: IDR functions (e.g., linker, entropic spring, PTM site clusters) are primarily determined by composition, with high tolerance for positional arrangement.
- Insertion/deletion of 5--20 residues in the same IDR typically does not affect function [Zibaee et al., \emph{Protein Sci.} 16, 906 (2007)].

\textbf{Situs's behavior on IDRs}:
\[I(Y; P_{\text{seq}} \mid S = \text{``IDR residue''}) \approx 0\]

By Theorems 2.2.1 and 2.3.1:
\[\delta_s^{\text{PE}} \leq \sqrt{\frac{1}{2} I(Y; P \mid S)} + \delta_s^{\text{variance}} \approx 0 + O(1/\sqrt{M})\]

That is: \textbf{Situs degenerates into additive noise of dimension $d$}.

Specifically, over $L$ consecutive positions in the IDR region:
\[\text{PE}(p) = [\sin(2\pi p/\lambda_1), \cos(2\pi p/\lambda_1), \ldots]\]

These encoding vectors vary smoothly with $p$ (Lipschitz continuity), but the variability they carry is \textbf{unrelated} to the label $Y$. During training:
- Clean samples: $\text{PE}(p)$ adds label-irrelevant fluctuations $\to$ may slightly reduce expert consistency $\to p_{\text{clean}}^{\text{PPE}} \geq p_{\text{clean}}$.
- Noisy samples: $\text{PE}(p)$ likewise provides no noise-discrimination ability $\to p_{\text{noisy}}^{\text{PPE}} \approx p_{\text{noisy}}$.
- Result: $\delta_s^{\text{PE}} \leq 0$ (negative or zero).

\textbf{Honest conclusion}: IDRs are Situs's \textbf{worst-case scenario} --- positional encoding degenerates into harmless but useless dimensional inflation. It does not actively damage the model (Lipschitz continuity guarantees $h_i$ remains smooth), but consumes $d$ dimensions of computation and storage, and may slightly reduce convergence speed through gradient noise. If the dataset has a high proportion of IDRs ($>$30\% as in eukaryotic proteins), Situs's net benefit \textbf{may be negative}.

\begin{center}\rule{0.5\linewidth}{0.5pt}\end{center}

\subsection{1.2 Materials 3D Structure Positional Encoding}\label{materials-3d-structure}

\subsubsection{1.2.1 $V_N$ Vacancy Formation Energy: Grain Boundary vs Bulk vs Surface}\label{vn-vacancy-formation}

\textbf{Physical facts (DFT computation data)}:

The formation energy of nitrogen vacancies $V_N$ in AlN depends strongly on spatial position [Freysoldt et al., \emph{Rev. Mod. Phys.} 86, 253 (2014); Stampfl \& Van de Walle, \emph{Phys. Rev. B} 65, 155212 (2002)]:

\begin{longtable}[]{@{}
  >{\raggedright\arraybackslash}p{(\linewidth - 4\tabcolsep) * \real{0.1714}}
  >{\raggedright\arraybackslash}p{(\linewidth - 4\tabcolsep) * \real{0.5429}}
  >{\raggedright\arraybackslash}p{(\linewidth - 4\tabcolsep) * \real{0.2857}}@{}}
\toprule\noalign{}
\begin{minipage}[b]{\linewidth}\raggedright
Position
\end{minipage} & \begin{minipage}[b]{\linewidth}\raggedright
Formation Energy $E_f$ (eV)
\end{minipage} & \begin{minipage}[b]{\linewidth}\raggedright
Physical Reason
\end{minipage} \\
\midrule\noalign{}
\endhead
\bottomrule\noalign{}
\endlastfoot
\textbf{Bulk} (coordination number 4) & 3.2--4.0 & Complete tetrahedral coordination, breaking 4 Al--N bonds \\
\textbf{Surface} ((0001) surface, coordination number 3) & 1.8--2.5 & Surface relaxation releases strain energy, reduced coordination \\
\textbf{Grain Boundary} ($\Sigma$7 symmetric tilt GB) & 1.0--1.8 & GB intrinsic stress field + excess volume + pre-existing dangling bonds \\
\end{longtable}

Difference magnitude: $\Delta E_f \approx 1.5$--$2.5$ eV, far exceeding $k_B T$ ($\approx$ 0.025 eV at room temperature), decisively dominant at \textbf{any} temperature.

\textbf{Physical causal chain} (irreducible):
\[\text{3D position} \xrightarrow{\text{determines}} \text{local coordination environment} \xrightarrow{\text{determines}} \text{defect formation energy}\]

This is a clear causal arrow --- position is the independent variable, formation energy is the dependent variable. Positional encoding captures the \textbf{cause side} of this causal chain.

\[I(Y = E_f; P_{\text{3D}} \mid S = V_N) > 0 \quad \text{--- strictly positive and large in the information-theoretic sense}\]

By the Fano lower bound (Theorem 2.4.1), for multi-class classification (different defect charge states):
\[\delta_s^{\text{PE}} \geq \frac{2 \cdot I(Y; P \mid S) - \log 2}{\log |\mathcal{Y}|}\]

Taking $I(Y; P \mid S) \approx H(Y|S) \approx \log 3 \approx 1.6$ bits (3 formation energy intervals: high/medium/low), $\log |\mathcal{Y}| \approx \log 3$:
\[\delta_s^{\text{PE}} \gtrsim 0.5 \text{--}0.8\]

This is a \textbf{substantial detection margin gain}, corresponding to significant tightening of Theorem 1's exponential bound.

\textbf{Honest conclusion}: For materials defect detection, 3D Situs is the \textbf{physically most valuable application scenario}. Position directly and causally determines the label. This is Situs's ``home turf.''

\begin{center}\rule{0.5\linewidth}{0.5pt}\end{center}

\subsubsection{1.2.2 Physical Symmetry Check of 3D Rotation Encoding}\label{3d-rotation-symmetry-check}

\textbf{Situs rotation encoding definition} (Definition 1.3.1):
\[\text{PE}_{\text{rot}}(\mathbf{p}) = \mathbf{R}_x(\alpha x) \cdot \mathbf{R}_y(\beta y) \cdot \mathbf{R}_z(\gamma z) \cdot \mathbf{e}_0\]

\textbf{Check 1: Translation invariance} \checkmark

Inner product:
\[\langle \text{PE}(\mathbf{p}), \text{PE}(\mathbf{q}) \rangle = \langle \mathbf{e}_0, \mathbf{R}(\mathbf{q} - \mathbf{p}) \mathbf{e}_0 \rangle = \cos(\alpha\Delta x) + \cos(\beta\Delta y) + \cos(\gamma\Delta z)\]

Depends only on $\Delta\mathbf{p} = \mathbf{q} - \mathbf{p}$. \textbf{Translation invariance holds strictly}. This satisfies the basic requirement of crystallography --- in a perfect crystal, two positions are equivalent if and only if their relative displacement is a lattice vector.

\textbf{Check 2: Rotation equivariance} $\times$ \textbf{Does not hold}

Consider a global rotation $R_{\text{phys}} \in SO(3)$ acting on coordinates: $\mathbf{p} \to R_{\text{phys}} \mathbf{p}$.

For $R_{\text{phys}} =$ 90$^\circ$ rotation about the $z$-axis (mapping $(x, y, z)$ to $(-y, x, z)$):
\[\text{PE}_{\text{rot}}(R_{\text{phys}}\mathbf{p}) = \mathbf{R}_x(-\alpha y) \mathbf{R}_y(\beta x) \mathbf{R}_z(\gamma z) \mathbf{e}_0\]

While $\text{PE}_{\text{rot}}(\mathbf{p}) = \mathbf{R}_x(\alpha x) \mathbf{R}_y(\beta y) \mathbf{R}_z(\gamma z) \mathbf{e}_0$.

There \textbf{does not exist} a $\mathbf{p}$-independent transformation $T$ such that $T \cdot \text{PE}(\mathbf{p}) = \text{PE}(R_{\text{phys}}\mathbf{p})$. Reason: The encoding assigns \textbf{disjoint dimension pairs} to the three coordinate axes --- $(0,1)$ for $x$, $(2,3)$ for $y$, $(4,5)$ for $z$. Physical rotation mixes $x$ and $y$ components, but there is no corresponding ``mixing'' mechanism in the encoding space.

\textbf{Consequences of the breakage}:

The absence of rotation equivariance is fatal in the following scenarios:
1. \textbf{Polycrystalline/grain boundary systems}: Different grains have different orientations; the encoding of the same type of defect differs due to orientation differences $\to$ $h_i$ introduces label-irrelevant variability.
2. \textbf{Molecular dynamics trajectories}: Overall rotation of the molecule causes all atom encodings to change $\to$ the model may learn ``rotational state'' rather than ``physical state.''
3. \textbf{Surface slab calculations}: Different cuts of the same surface (e.g., AlN (0001) vs (11$\bar{0}$0)) have different coordinate axis orientations.

\textbf{Can it be fixed?}

Theoretically yes. Alternatives:
- \textbf{Spherical harmonic encoding}: $\text{PE}(\mathbf{p}) = [Y_{\ell m}(\theta, \phi)]$, with strict $SO(3)$ rotation equivariance.
- \textbf{E(3) equivariant networks}: Using Tensor Field Networks [Thomas et al., \emph{NeurIPS} 2018] or SE(3)-Transformers [Fuchs et al., \emph{NeurIPS} 2020], naturally possessing translation + rotation equivariance.
- \textbf{Cost}: The dimension of spherical harmonic encoding grows as $(L+1)^2$ with angular momentum $L$; $L=3$ already requires 16 dimensions (vs 6 dimensions for rotation encoding), and the Lipschitz behavior is more complex.

\textbf{Honest conclusion}: Situs's 3D rotation encoding \textbf{only possesses translation invariance, not rotation equivariance}. For most materials applications (requiring handling of arbitrarily oriented crystals), this is a \textbf{serious defect}. However, in controlled scenarios (single crystals, fixed-orientation slabs, orientation-aligned datasets), this is not a problem.

\begin{center}\rule{0.5\linewidth}{0.5pt}\end{center}

\subsection{1.3 Total Atom Count $N$: DFT Accuracy for 32 vs 256 Supercells}\label{atom-count-dft-accuracy}

\subsubsection{1.3.1 Physical Origin of Systematic Differences}\label{systematic-differences-physical-origin}

\textbf{Makov-Payne finite-size scaling} [Makov \& Payne, \emph{Phys. Rev. B} 51, 4014 (1995)]:

The scaling behavior of charged defect formation energy with supercell size $L$:
\[E_f(L) = E_f(\infty) + \frac{\alpha q^2}{2\varepsilon L} + \frac{2\pi q Q}{3\varepsilon L^3} + O(L^{-5})\]

where $q$ is the defect charge, $\varepsilon$ is the dielectric constant, $Q$ is the elastic quadrupole moment.

\textbf{Specific numerical estimates ($V_N^{+1}$ in AlN, $\varepsilon \approx 9$)}:

\begin{longtable}[]{@{}
  >{\raggedright\arraybackslash}p{(\linewidth - 10\tabcolsep) * \real{0.1538}}
  >{\raggedright\arraybackslash}p{(\linewidth - 10\tabcolsep) * \real{0.1231}}
  >{\raggedright\arraybackslash}p{(\linewidth - 10\tabcolsep) * \real{0.2308}}
  >{\raggedright\arraybackslash}p{(\linewidth - 10\tabcolsep) * \real{0.1692}}
  >{\raggedright\arraybackslash}p{(\linewidth - 10\tabcolsep) * \real{0.2000}}
  >{\raggedright\arraybackslash}p{(\linewidth - 10\tabcolsep) * \real{0.1231}}@{}}
\toprule\noalign{}
\begin{minipage}[b]{\linewidth}\raggedright
Supercell Size
\end{minipage} & \begin{minipage}[b]{\linewidth}\raggedright
Atom Count
\end{minipage} & \begin{minipage}[b]{\linewidth}\raggedright
Periodic Image Separation
\end{minipage} & \begin{minipage}[b]{\linewidth}\raggedright
$1/L$ Correction
\end{minipage} & \begin{minipage}[b]{\linewidth}\raggedright
$1/L^3$ Correction
\end{minipage} & \begin{minipage}[b]{\linewidth}\raggedright
Total Correction
\end{minipage} \\
\midrule\noalign{}
\endhead
\bottomrule\noalign{}
\endlastfoot
2$\times$2$\times$2 & 32 & $\sim$6.2 \AA & $\sim$0.35 eV & $\sim$0.12 eV & \textbf{$\sim$0.5 eV} \\
2$\times$2$\times$3 & 48 & $\sim$6.2--9.3 \AA & $\sim$0.25 eV & $\sim$0.05 eV & \textbf{$\sim$0.3 eV} \\
3$\times$3$\times$3 & 108 & $\sim$9.3 \AA & $\sim$0.15 eV & $\sim$0.02 eV & \textbf{$\sim$0.17 eV} \\
4$\times$4$\times$4 & 256 & $\sim$12.4 \AA & $\sim$0.08 eV & $\sim$0.005 eV & \textbf{$\sim$0.09 eV} \\
\end{longtable}

\textbf{Physical origins} (three independent finite-size effects, all non-negligible):

\begin{enumerate}
\def\labelenumi{\arabic{enumi}.}
\tightlist
\item
  \textbf{Electrostatic image charge interaction} ($1/L$ dominant term):
  \begin{itemize}
  \tightlist
  \item
    Coulomb interaction between the charged defect and its periodic images.
  \item
    In a 32-atom supercell, adjacent images are only 6.2 \AA\ apart $\to$ interaction energy $\sim$0.3--0.5 eV.
  \item
    Physically this is artificially introduced --- a real isolated defect has no periodic images.
  \end{itemize}
\item
  \textbf{Elastic image interaction} ($1/L^3$):
  \begin{itemize}
  \tightlist
  \item
    Interaction between the local strain field induced by the defect and its periodic images.
  \item
    For $V_N$ (removing an atom produces tensile strain), this effect is non-negligible.
  \end{itemize}
\item
  \textbf{Defect wavefunction overlap} (exponential decay $\propto e^{-L/\xi}$):
  \begin{itemize}
  \tightlist
  \item
    For deep-level defects (e.g., $V_N$ in AlN produces a $\sim$1.5 eV deep donor level), the wavefunction localization radius $\xi \approx 2$--$3$ \AA.
  \item
    In a 32-atom supercell, the overlap of adjacent wavefunction tails is non-negligible $\to$ defect level splits into a band (artificial dispersion).
  \end{itemize}
\end{enumerate}

\textbf{What does this have to do with Situs's $N$ encoding?}

Situs's third domain (Definition 1.1.3) encodes the total atom count $N$ as a scalar position. The problem is:

\begin{itemize}
\tightlist
\item
  $N=32$ and $N=256$ data are \textbf{not different ``positions'' of the same physical system} --- they are approximations at \textbf{different accuracy levels}.
\item
  $N$ encoding confuses ``computational accuracy'' with ``physical label.'' The $E_f$ deviation at $N=32$ is a computational artifact, not a physical signal.
\item
  If the training set mixes data from 32- and 256-atom supercells, $N$ encoding may learn the spurious correlation ``$N=32 \to$ high $E_f$.''
\end{itemize}

\textbf{Physically correct treatment}:
- Apply Freysoldt-Neugebauer-Van de Walle (FNV) finite-size correction [Freysoldt et al., \emph{Phys. Rev. Lett.} 102, 016402 (2009)] to $E_f$, eliminating the systematic $N$-dependent bias.
- \textbf{Then} consider whether to use $N$ encoding --- at this point $N$ reflects genuine system size effects (e.g., quantum confinement), not computational artifacts.

\textbf{Honest conclusion}: If the training data has not undergone finite-size correction, $N$ encoding learns \textbf{computational artifacts rather than physics}. After correction, the genuine physical information carried by $N$ (quantum confinement, number of vibrational modes) is small for defect physics ($\Delta E_f$'s $N$-dependence $<$ 0.05 eV for $N > 100$). For most materials tasks, $N$ encoding's $\delta_s^{\text{PE}}$ is \textbf{close to zero} --- this is a physically weakly motivated, practically dangerous (prone to learning artifacts) encoding.

\begin{center}\rule{0.5\linewidth}{0.5pt}\end{center}

\section{Part 2: Applicability Scenario Summary}\label{part-2-applicability-scenario-summary}

\begin{longtable}[]{@{}
  >{\raggedright\arraybackslash}p{(\linewidth - 6\tabcolsep) * \real{0.1667}}
  >{\raggedright\arraybackslash}p{(\linewidth - 6\tabcolsep) * \real{0.1667}}
  >{\raggedright\arraybackslash}p{(\linewidth - 6\tabcolsep) * \real{0.3333}}
  >{\raggedright\arraybackslash}p{(\linewidth - 6\tabcolsep) * \real{0.3333}}@{}}
\toprule\noalign{}
\begin{minipage}[b]{\linewidth}\raggedright
Scenario
\end{minipage} & \begin{minipage}[b]{\linewidth}\raggedright
$I(Y;P|S) > 0$?
\end{minipage} & \begin{minipage}[b]{\linewidth}\raggedright
Situs Benefit
\end{minipage} & \begin{minipage}[b]{\linewidth}\raggedright
Key Limitation
\end{minipage} \\
\midrule\noalign{}
\endhead
\bottomrule\noalign{}
\endlastfoot
Protein active site vs surface (3D) & Strong YES & \textbf{High} & Requires 3D structure; 1D encoding weak \\
$\alpha$-helix periodicity & Weak YES & \textbf{Low-Medium} & $d < 167$ cannot resolve; explicit period features needed \\
IDR regions & NO & \textbf{None or Negative} & Encoding degenerates to noise; wastes dimensions \\
Materials defect formation energy (3D) & Strong YES & \textbf{High} & Home turf; position $\to$ coordination $\to$ energy \\
Polycrystalline / arbitrary orientation & YES but orientation-mixed & \textbf{Low} & Lack of rotation equivariance introduces artifacts \\
Supercell size $N$ (uncorrected) & NO (artifact) & \textbf{Negative} & Learns computational artifacts; apply FNV correction first \\
Supercell size $N$ (corrected) & Weak YES & \textbf{Low} & Genuine size effects small for most properties \\
\end{longtable}

\begin{center}\rule{0.5\linewidth}{0.5pt}\end{center}

\emph{End of situs\_physical\_validation.tex}

\end{document}
